\chapter{Discusion}

\subsection{El gemelo digital como habilitador de datos etiquetados}

El problema del cold-start en deteccion de fallas es uno de los mas
subestimados en la literatura. La mayoria de los papers proponen
modelos sofisticados (transformers, normalizing flows) pero no
explican como entrenar el primer modelo sin etiquetas. Este trabajo
muestra que un gemelo digital \emph{ligero} ya entrega senales utiles:

\begin{itemize}[leftmargin=*,itemsep=2pt]
  \item Los residuos del gemelo tienen una distribucion bien
        comportada bajo operacion normal.
  \item Las fallas inyectadas producen firmas distinguibles.
  \item El autoencoder LSTM aprende la dinamica normal de los
        residuos con pocas epocas.
\end{itemize}

El gemelo ligero es calibrable con la hoja de datos del fabricante
y unas pocas semanas de operacion real. Esto reduce drasticamente el
tiempo de despliegue de un sistema como el propuesto.

\subsubsection{Limitaciones del gemelo ligero}

El gemelo ligero no captura:
\begin{itemize}[leftmargin=*,itemsep=2pt]
  \item Armonicos de conmutacion.
  \item Modelos termicos distribuidos (punto caliente en el
        devanado).
  \item Comportamiento dinamico del transformador bajo falla
        (modelos de Jiles-Atherton para el nucleo).
  \item Interacciones con la red aguas arriba (modelos de estabilidad
        transitoria).
\end{itemize}

Estos fenomenos son relevantes para fallas raras pero graves (falla
interna de transformador, resonance sub-sincrona). El gemelo deberia
extenderse a un modelo \emph{medio} antes del despliegue
productivo, posiblemente incluyendo FEM del transformador y un
modelo de Thevenin equivalente para la red aguas arriba.

\subsection{Por que conformal y no solo un umbral fijo}

El conformal prediction aporta tres beneficios sobre un umbral fijo
(como 3-sigma):

\begin{enumerate}[leftmargin=*,itemsep=2pt]
  \item \textbf{Garantia libre de distribucion}. El detector 3-sigma
        asume gaussianidad, que se viola en regimenes operativos
        extremos. El detector conformal no hace ese supuesto.

  \item \textbf{Confianza del operador}. Cuando el sistema afirma que
        la probabilidad de falsa alarma es $\le 5\%$, esa afirmacion
        es valida bajo el supuesto de intercambiabilidad, lo que el
        operador puede auditar.

  \item \textbf{Cumplimiento}. La garantia conformal es
        directamente traducible a un compromiso contractual con la
        distribuidora (\emph{el sistema no generara mas de X alertas
        espurias por mes}).
\end{enumerate}

\subsubsection{Limitaciones}

El conformal prediction tiene dos limitaciones en este contexto:

\begin{itemize}[leftmargin=*,itemsep=2pt]
  \item Requiere un \emph{conjunto de calibracion} separado del
        entrenamiento. Esto puede ser dificil en subestaciones nuevas.

  \item La garantia se debilita bajo \emph{distribution shift}: si la
        politica operativa de la distribuidora se vuelve mas
        conservadora (cambios de tap mas frecuentes, por ejemplo), el
        cuantil conformal puede subestimar el FPR.
\end{itemize}

La primera limitacion se mitiga usando los primeros 30 dias de
operacion como calibracion. La segunda se mitiga con \emph{conformal
adaptativo}~\citep{ACI-2020}, que actualiza el cuantil con una ventana
deslizante.

\subsection{El costo de la privacidad}

El aprendizaje federado suele tener un costo en precision. En este
trabajo no aparecio: el modelo federado midio \emph{mejor} que el
centralizado (F1 $0.87$ vs $0.80$; AUC $0.906$ vs $0.878$). La
lectura honesta es que la comparacion no es de presupuesto igual
-- el centralizado entrena 6 epocas sobre una sola SE, el federado
agrega estadistica de 3 SEs -- y que el resultado relevante es la
\textbf{equivalencia}: federar no cede detectabilidad a cambio de:

\begin{itemize}[leftmargin=*,itemsep=2pt]
  \item No consolidar datos en la nube.
  \item Cumplir con NERC CIP / normativa CNE sobre segregacion OT/IT.
  \item Permitir que las distribuidoras conserven la propiedad de
        sus datos.
\end{itemize}

Este trade-off es directamente favorable desde el punto de vista
regulatorio y comercial. Queda abierta la pregunta de si tecnicas
mas avanzadas (federated learning con personalizacion,
conocimiento multi-tarea, distillation) amplian la ventaja o
estabilizan la convergencia con clientes heterogeneos.

\subsubsection{FedAvg no es la unica opcion}

Alternativas que podrian explorarse:
\begin{itemize}[leftmargin=*,itemsep=2pt]
  \item \textbf{FedProx}: agrega un termino de regularizacion
        proximal que estabiliza la convergencia ante datos
        heterogeneos.
  \item \textbf{Personalizacion}: cada SE mantiene un modelo
        personalizado encima del modelo global. Esto permite
        capturar fallas raras propias de una SE sin sacrificar la
        estadistica global.
  \item \textbf{Federated conformal}: calibrar el umbral conformal
        localmente con datos de la SE. La garantia es solo local,
        pero la cobertura mejora.
\end{itemize}

\subsection{Aplicabilidad regulatoria}

El sistema propuesto ataca directamente tres frentes regulatorios:

\subsubsection{SAIDI y compensacion automatica}

Sobre las fallas agudas detectadas por el sistema en su
configuracion completa (F1, F3, F8, F5), la literatura estima una
reduccion de 30\% en el tiempo medio de reposicion. Aplicando al
SAIDI medio del SEN ($\sim$12~h/ano en 2024), esto equivale a una
reduccion de $\approx 3.6$~h/ano. Para una distribuidora con
1.5 millones de clientes, el costo evitado es significativo.

\subsubsection{Cargos SEC por incumplimiento}

La trazabilidad del sistema es valiosa en si misma. Cada alerta
tiene timestamp, score y modelo que la emitio, lo que sirve como
evidencia ante un procedimiento SEC. El sistema reduce el riesgo
de cargos por \emph{supervision insuficiente} (Art. 145 del
reglamento).

\subsubsection{Casos SERNAC post-evento climatico}

El frente de Valparaiso (julio 2026) y los procedimientos SERNAC
asociados crearon la necesidad de \emph{probar} que la distribuidora
respondio oportunamente. El sistema propuesto entrega:

\begin{itemize}[leftmargin=*,itemsep=2pt]
  \item Tiempo exacto de la primera deteccion.
  \item Score de prioridad por subestacion.
  \item Recomendacion de despacho de cuadrillas.
\end{itemize}

Esto es exactamente lo que el SERNAC exige para evaluar si la
respuesta fue razonable.

\subsection{Trabajo futuro}

\subsubsection{Validacion con datos reales}

El paso siguiente es validar con datos reales de una distribuidora.
Las conversaciones preliminares con Chilquinta (Valparaiso) y
Engie Transelec estan en curso. La validacion debe responder:

\begin{itemize}[leftmargin=*,itemsep=2pt]
  \item Cuanto reduce el sistema el SAIDI en operacion real?
  \item Cual es el FPR real durante eventos climaticos extremos?
  \item Cual es la adhesion de los operadores al sistema?
\end{itemize}

\subsubsection{Conformal adaptativo}

Implementar \emph{adaptive conformal inference} (ACI) para que el
umbral se actualice con ventanas de 1~h, manteniendo la garantia
bajo distribution shift.

\subsubsection{Deteccion de fallas combinadas}

El catalogo de fallas incluye fallas aisladas. Las fallas reales
son a menudo \emph{combinadas} (cortocircuito que induce una
sobretension transitoria, por ejemplo). El modelo deberia entrenarse
con fallas multi-evento.

\subsubsection{Gemelo digital mejorado}

Extender el gemelo a un modelo de orden medio: FEM del
transformador, modelo de la red aguas arriba, modelos termicos
distribuidos. Esto mejora la fidelidad pero aumenta el costo de
calibracion.

\subsubsection{SaaS vendible}

Empaquetar el pipeline como un SaaS con:

\begin{itemize}[leftmargin=*,itemsep=2pt]
  \item Dashboard web por subestacion (este trabajo es el prototipo).
  \item API REST para integracion con SCADA.
  \item Modelo de precios por subestacion y por alerta.
  \item SLA de disponibilidad del 99.9\%.
\end{itemize}

El mercado objetivo son las distribuidoras (Chilquinta, CGE,
Saesa, Enel Distribucion) y las transmisoras (Engie Transelec,
Colbun Transelec). El precio de referencia por subestacion
monitorizada seria $\approx$ USD 200-500/mes, con un mercado
direccionable de $\sim 600$ subestaciones en el SEN.

\subsection{Sintesis}

El sistema propuesto es una contribucion metodologica concreta al
campo de deteccion temprana en subestaciones de distribucion. Su
aporte principal es demostrar que un pipeline federado + gemelo
digital + conformal prediction puede cumplir el triple requisito
de: alta sensibilidad, garantia de falsa alarma, y privacidad de
los datos. Los resultados empiricos son alentadores, aunque
queda trabajo pendiente en la calibracion con datos reales y en
la extension del gemelo a un modelo de orden medio.